Considere $X=\{x_i\}_{i\in I}$ conjunto de asociamos e $X^{-1}=\{x_i^{-1}\}$, asociamos a $X$ seu conjugado.\\
\begin{definition}[Alfabeto e palabras]
  Consider $X\cup X^{-1}$ chamado \textbf{alfabeto}. Assim, construimos \textbf{palabras} com sequencias finitas de elementos de $X\cup X^{-1}$ como $x_{i_{1}}^{\pm}x_{i_{2}}^{\pm}\cdots x_{i_{n}}^{\pm}$.\\
  Chamamos $x_ix_i^{-1}$ \textbf{palabra vazia}.\\
  Mas vamos considerar o conjunto $F(X)$ como o conjunto de so as palabras reduzidas, isto é, que elas não tem palabras vazias.
\end{definition}
Sendo así, consideramos uma justaposição de palabras como operação e exigendo como resultado uma palabra reduzida.
\begin{example}
  Seja $v=x_5x_4x_1^{-1}x_2x_3^{-1}$ e $w=x_3x_2^{-1}x_3x_1$, logo $vw=x_5x_4x_1^{-1}x_3x_1$. 
\end{example}
\begin{definition}[Grupo livre]
  $F(X)$ é um grupo, dito \textbf{grupo livre} gerado por $X$. 
\end{definition}
Note que $|X|$ é posto de grupo livre $F(X)$.
\begin{note}
  Existe um grupo livre de qualquer posto dado (finito ou não).
\end{note}
Note que 
\begin{itemize}
  \item $\mathbb{Z}$ é um grupo livre com posto $1$.
  \item Um grupo livre $F(X)$ é infinito. 
\end{itemize}
\textbf{Questão}: dado um grupo infinito, temos
\begin{enumerate}
  \item Se $o(x)$ é finito, para todo $x\in G$, então $G$ não comtém grupos livres.
  \item Se existe $x\in G$ tal que $o(x)$ é infinito, então $G$ comtém um grupo livre de posto $1$. 
\end{enumerate}
\textbf{Pregunta}: $G$ comtém grupos livres de posto $\geq 2$?
\begin{theorem}
  Se $F$ é um corpo infinito, então o grupo $GL_{n}(F)$ comtém subgrupos livres de posto $\geq 2$, para $n\geq 2$. 
\end{theorem}
\begin{lemma}[Macbeath]
  Se $G$ é um grupo gerado por subgrupos $G_1$ e $G_2$ com $G_1\cap G_2=\{1\}$ tais que $|G_i|> 2$, para algúm $i$ e $G$ ace em um conjunto com ações.\\
  $\Omega=\Omega_1\cup\Omega_2$ unión disjunta.
  \begin{itemize}
    \item para todo $g\in G_1\setminus\{1\}$ temos que $g_1(\Omega_1)\subseteq\Omega_{2}$.
    \item para todo $g\in G_2\setminus\{1\}$ temos que $g_2(\Omega_{2})\subseteq\Omega_1$. 
  \end{itemize}
  Então, temos que $G$ é o produto livre de $G_1$ e $G_2$. 
\end{lemma}
\begin{proof} 
  Vamos a provar que não existe uma relação não trivial do tipo
  \begin{align*}
    x_1y_1x_2y_2\cdots=1,
  \end{align*}
  com $x_i\in G_r\setminus\{1\},y_i\in G_s\setminus\{1\},\,r,s\in\{1,2\}$ e $r\neq s$.\\
  Raciocinemos pelo contradição.\\
  Suponha que $|G_1|\geq 3$ e definamos $w$ a palavra dada por
  \begin{align*}
    w=x_1y_1x_2y_2\cdots=1.
  \end{align*}
  Sem perdida de generalidade nos podemos asumir que $w$ começa e termina com um elemento do $G_1$. Pelo contrario, se $w$ començã com un elemento de $G_1$ e termina com um elementto de $G_2$ podemos fazer o seguinte, suponha que $g\in G_1$ com $g\notin \{1,x_1\}$, logo temos que
  \begin{align*}
    g_1^{-1}wg_1=g_1^{-1}1g_1=1.
  \end{align*}
  Isto é
  \begin{align*}
    w=(g_1^{-1}x_1)y_1x_2y_2\cdots x_ny_ng_1=1.
  \end{align*}
  Portanto, podemos afirmar que $w$ comença y termina com um elemento de $G_1$, assim temos que
  \begin{align*}
    w=x_1y_1x_2y_2\cdots x_ny_nx_{n+1}=1.
  \end{align*}
  Então, como $w=1$, temos que $w(\Omega_1)=\Omega_1$. Por outro, pela hipótese temos que
  \begin{align*}
    x_1y_1x_2y_2\cdots x_ny_nx_{n+1}(\Omega 1)&\subseteq x_1y_1x_2y_2\cdots x_ny_n(\Omega_2)\\
    &\subseteq x_1y_1x_2y_2\cdots x_n(\Omega_1)\\
    &\subseteq x_1y_1x_2y_2\cdots (\Omega_2)\\
    &\subseteq x_1y_1x_2y_2(\Omega_2)\\
    &\subseteq x_1y_1x_2(\Omega_1)\\
    &\subseteq x_1y_1(\Omega_2)\\
    &\subseteq x_1(\Omega_1)\\
    &\subseteq (\Omega_2).
  \end{align*}
  Isto é
  \begin{align*}
    w(\Omega_1)\subseteq \Omega_2.
  \end{align*}
  Mas $w(\Omega_1)=\Omega_1$, pelo que temos que $\Omega_1\subseteq \Omega_2$. Absurdo, ya que $\Omega_1\cup\Omega_2$ são uma união disjunta. Assim, temos que não existe uma relação não trivial do tipo
  \begin{align*}
    w=x_1y_1x_2y_2\cdots=1.
  \end{align*}
  Logo como toda palavra reducida e sem elemento neutro é distinta de $1$, temos que $G$ é o produto livre de $G_1\ast G_2$.
\end{proof}
\begin{example}[grupo livre de permutações do plano complexo]
  Seja $\alpha\in \mathbb{C}$ com $|\alpha|\geq 2$. Considere $f$ e $g$ dadas por lo siguiente
  \begin{align*}
    f(z)&=z+\alpha.\\
    g(z)&=\begin{cases}
      \frac{z}{\alpha z+1}, &\text{ se } z\neq-\frac{1}{\alpha}\text{,} \\
      \frac{1}{\alpha}, &\text{ se } z=-\frac{1}{\alpha}.
    \end{cases}
  \end{align*}
  Veamos que $\left\langle f,g \right\rangle$ é um grupo livre.\\
  Primeiro, note que para qualquer $z\in \mathbb{C}$, em particular com $|z|< 1$ temos que
  \begin{align*}
    |z+\alpha| \geq |\alpha|-|z| > 2-1=1.
  \end{align*}
  Mais geralmente temos que se $k\geq 1$ temos que 
  \begin{align*}
    |f^{k}(z)|=|z+k\alpha|\geq k|\alpha|-|z|> 2k-1\geq 1.
  \end{align*}
  Logo, podemos dizer que $f$ envia o interior do círculo unitário para o exterior do círculo unitário.\\
  Agora, veamos que $g$ envia o exterior do círculo unitário para o interior do círculo unitário, ja que se pegamos $z\in\mathbb{C}$ com $|z|< 1$ temos que 
  \begin{align*}
    g\left(\frac{1}{z} \right)=\frac{\frac{1}{z}}{\alpha\frac{1}{z}+1}=\frac{1}{\alpha+z}.
  \end{align*}
  Mais geralmente, temos que se $k\geq 1$ temos que
  \begin{align*}
    g^{k}\left( \frac{1}{z} \right)=\frac{1}{k\alpha+z}.
  \end{align*}
  logo como temos que
  \begin{align*}
    |k\alpha+z|\geq |k\alpha|-|z|>2k-1\geq 1.
  \end{align*}
  Assim,
  \begin{align*}
    \left| g^{k}\left( \frac{1}{z} \right) \right| = \left| \frac{1}{k\alpha+z} \right|< 1.
  \end{align*}
  Assim, nos definimos
  \begin{align*}
    \Omega_1&=\{z\in\mathbb{C}: |z|< 1\}.\\
    \Omega_2&=\{z\in\mathbb{C}: |z|> 1\}=\left\{\frac{1}{w}\in\mathbb{C}:|w|<1\right\}.
  \end{align*}
  Assim, temos que $\left\langle g \right\rangle\cap \left\langle f \right\rangle=\{1\}$ com ambos de ordem mais grande que $2$ e com ações mencionadas anteriormente para $\Omega_1$ e $\Omega_2$ que cumplen $f^{k}(\Omega_1)\subseteq \Omega_2$ e $g^{k}(\Omega_2)\subseteq \Omega_1$ para todo $k$, logo temos que pelo lema de Macbeath $\left\langle f,g \right\rangle$ é um grupo livre. 
\end{example}
\begin{theorem}[Sanov]
  Considere $\alpha\in\mathbb{C}$, $|\alpha|\geq 2$. Então
  \begin{align*}
    A=\begin{pmatrix}
      1 & \alpha \\
      0 & 1
    \end{pmatrix}\quad B=\begin{pmatrix}
      1 & 0 \\
      \alpha & 0
    \end{pmatrix}.
  \end{align*}
  geram um subgrupo livre de $GL_2(\mathbb{C})$.
\end{theorem}
\begin{proof}
  Considerando $a,b,c,d\in\mathbb{C}$ com $ad-bd\neq 0$, definimos:
  \begin{align*}
    \lambda_{a,b,c,d}:\mathbb{C}&\to\mathbb{C},\\
    \xi &\to \begin{cases}
      \frac{a\xi+b}{c\xi+d}, &\text{ para } \xi\neq-\frac{d}{c} \text{,} \\
      \frac{d}{c}, &\text{ para } \xi=-\frac{d}{c}.
    \end{cases}
  \end{align*}
  Com isso, dada 
  \begin{align*}
    \begin{pmatrix}
      a & b \\
      c & d
    \end{pmatrix}\in GL_{2}(\mathbb{C})
  \end{align*}
  podemos definir
  \begin{align*}
    \varphi:GL_{2}(\mathbb{C})&\to\mathcal{G},\\
    \begin{pmatrix}
      a & b \\
      c & d
    \end{pmatrix}&\to \lambda_{a,b,c,d},
  \end{align*}
  onde $\mathcal{G}=\{\lambda_{a,b,c,d}:\mathbb{C}\to\mathbb{C}:ab-bd\neq 0\}$.\\
  Note que $\varphi$ é um homomorfismo sobrejetor de grupos, assim
  \begin{align*}
    \varphi(A)=\lambda_{1,\alpha,0,1},\quad f=\lambda_{1,\alpha,0,1}(\xi)=\xi+a.\\
    \varphi(B)=\lambda_{1,0,\alpha,1},\quad g=\lambda_{1,0,\alpha,1}(\xi)=\begin{cases}
      \frac{\xi}{\alpha\xi+1}, &\text{ para } \xi\neq-\frac{-1}{\alpha} \text{,} \\
      \frac{1}{\alpha}, &\text{ com } \xi=-\frac{1}{\alpha}.
    \end{cases}
  \end{align*}
  Assim, temos que pelo exemplo anterior $\left\langle f,g \right\rangle$ é um grupo livre e isto garante que $\left\langle A,B \right\rangle$ é um grupo livre.
\end{proof}
\begin{example}
  Seja $G_1=S_3=\{1,(12),(13),(23),(123),(132)\}$ e\\ $G_2=D_{4}=\{1,r,\theta,\theta^{2},\theta^{3},r\theta,r\theta^{2},r\theta^{3}\}$.\\
  Logo $G_1,G_2\leq G_1*G_2$ com $g_{i_1}g_{i_2}\cdots g_{i_k}$, onde $g_{i_j}\in G_1\cup G_2$
  \begin{align*}
    (12)\theta(13)(12)r\theta^{2}\theta^2=(12)\theta(123)r.
  \end{align*}
\end{example}
O grupo livre de $G_1$ e $G_2$ é formado por palavras reduzidas que são sequencias de elementos em $G_1$ e $G_2$, respetando as relações entre elementos individualizados de $G_1$ e $G_2$ respetando sem relação entre os elementos de $G_1$ e $G_2$.\\
Consider $G_1=\left\langle X,R \right\rangle$ e $G_2=\left\langle Y,S \right\rangle$ onde $X$ e $Y$ são conjuntos geradores e $R$ e $S$ são relações, então $G_1*G_2=\left\langle X,Y,R,S \right\rangle$.\\
Se $G_1$ e $G_2$ são grupos livres então, temos que $G_1=\left\langle X,\emptyset \right\rangle=\left\langle X|\phantom{\emptyset} \right\rangle$ e $G_2=\left\langle Y,\emptyset \right\rangle=\left\langle Y|\phantom{\emptyset} \right\rangle$, logo $G_1*G_2=\left\langle X,Y|\phantom{\emptyset} \right\rangle$.
\begin{example}
  De novo, $G_1=S_3$ e $G_2=D_4$. 
\end{example}
